% problem_id: S001-proposition-length-determinant-periodic-complex
% source_id: S001 (Stacks Project)
% source_file: chow.tex
% statement_locator: chow.tex lines 18023--18046
% proof_locator: chow.tex lines 18048--18272
% env: proposition
% status: text-extracted-verbatim (difficulty judgement pending, written by hand)
%
% The statement and proof below are the author's own text, extracted verbatim.
% Nothing is paraphrased, shortened, or supplemented.

\begin{proposition}
\label{proposition-length-determinant-periodic-complex}
Let $R$ be a local Noetherian ring with residue field $\kappa$.
Suppose that $(M, \varphi, \psi)$ is a $(2, 1)$-periodic
complex over $R$. Assume
\begin{enumerate}
\item $M$ is a finite $R$-module,
\item the cohomology modules of $(M, \varphi, \psi)$ are of finite length, and
\item $\dim(\text{Supp}(M)) = 1$.
\end{enumerate}
Let $\mathfrak q_i$, $i = 1, \ldots, t$ be the minimal
primes of the support of $M$. Then we have\footnote{
Obviously we could get rid of the minus sign by redefining
$\det_\kappa(M, \varphi, \psi)$ as the inverse of its
current value, see Definition \ref{definition-periodic-determinant}.}
$$
- e_R(M, \varphi, \psi) =
\sum\nolimits_{i = 1, \ldots, t}
\text{ord}_{R/\mathfrak q_i}\left(
\det\nolimits_{\kappa(\mathfrak q_i)}
(M_{\mathfrak q_i}, \varphi_{\mathfrak q_i}, \psi_{\mathfrak q_i})
\right)
$$
\end{proposition}

\begin{proof}
We first reduce to the case $t = 1$ in the following way.
Note that
$\text{Supp}(M) = \{\mathfrak m, \mathfrak q_1, \ldots, \mathfrak q_t\}$,
where $\mathfrak m \subset R$ is the maximal ideal.
Let $M_i$ denote the image of $M \to M_{\mathfrak q_i}$,
so $\text{Supp}(M_i) = \{\mathfrak m, \mathfrak q_i\}$.
The map $\varphi$ (resp.\ $\psi$) induces an $R$-module map
$\varphi_i : M_i \to M_i$ (resp.\ $\psi_i : M_i \to M_i$).
Thus we get a morphism of $(2, 1)$-periodic complexes
$$
(M, \varphi, \psi) \longrightarrow
\bigoplus\nolimits_{i = 1, \ldots, t} (M_i, \varphi_i, \psi_i).
$$
The kernel and cokernel of this map have support contained in
$\{\mathfrak m\}$. Hence by Lemma \ref{lemma-compare-periodic-lengths}
we have
$$
e_R(M, \varphi, \psi) =
\sum\nolimits_{i = 1, \ldots, t}
e_R(M_i, \varphi_i, \psi_i)
$$
On the other hand we clearly have $M_{\mathfrak q_i} = M_{i, \mathfrak q_i}$,
and hence the terms of the right hand side of the formula of the
lemma are equal to the expressions
$$
\text{ord}_{R/\mathfrak q_i}\left(
\det\nolimits_{\kappa(\mathfrak q_i)}
(M_{i, \mathfrak q_i}, \varphi_{i, \mathfrak q_i}, \psi_{i, \mathfrak q_i})
\right)
$$
In other words, if we can prove the lemma for each of the modules
$M_i$, then the lemma holds. This reduces us to the case $t = 1$.

\medskip\noindent
Assume we have a $(2, 1)$-periodic complex $(M, \varphi, \psi)$
over a Noetherian local ring with $M$ a finite $R$-module,
$\text{Supp}(M) = \{\mathfrak m, \mathfrak q\}$, and
finite length cohomology modules. The proof in this case
follows from Lemma \ref{lemma-key-lemma} and careful bookkeeping.
Denote
$K_\varphi = \Ker(\varphi)$,
$I_\varphi = \Im(\varphi)$,
$K_\psi = \Ker(\psi)$, and
$I_\psi = \Im(\psi)$.
Since $R$ is Noetherian these are all finite $R$-modules.
Set
$$
a = \text{length}_{R_{\mathfrak q}}(I_{\varphi, \mathfrak q})
= \text{length}_{R_{\mathfrak q}}(K_{\psi, \mathfrak q}),
\quad
b = \text{length}_{R_{\mathfrak q}}(I_{\psi, \mathfrak q})
= \text{length}_{R_{\mathfrak q}}(K_{\varphi, \mathfrak q}).
$$
Equalities because the complex becomes exact after localizing at
$\mathfrak q$. Note that $l = \text{length}_{R_{\mathfrak q}}(M_{\mathfrak q})$
is equal to $l = a + b$.

\medskip\noindent
We are going to use Lemma \ref{lemma-good-sequence-exists}
to choose sequences of elements in finite $R$-modules
$N$ with support contained in $\{\mathfrak m, \mathfrak q\}$.
In this case $N_{\mathfrak q}$ has finite length, say $n \in \mathbf{N}$.
Let us call a sequence $w_1, \ldots, w_n \in N$
with properties (1) and (2) of Lemma \ref{lemma-good-sequence-exists}
a ``good sequence''. Note that the quotient
$N/\langle w_1, \ldots, w_n \rangle$ of $N$ by the submodule generated by
a good sequence has support (contained in) $\{\mathfrak m\}$
and hence has finite length (Algebra, Lemma \ref{algebra-lemma-support-point}).
Moreover, the symbol
$[w_1, \ldots, w_n] \in \det_{\kappa(\mathfrak q)}(N_{\mathfrak q})$
is a generator, see Lemma \ref{lemma-determinant-dimension-one}.

\medskip\noindent
Having said this we choose good sequences
$$
\begin{matrix}
x_1, \ldots, x_b & \text{in} & K_\varphi, &
t_1, \ldots, t_a & \text{in} & K_\psi, \\
y_1, \ldots, y_a & \text{in} & I_\varphi \cap \langle t_1, \ldots t_a\rangle, &
s_1, \ldots, s_b & \text{in} & I_\psi \cap \langle x_1, \ldots, x_b\rangle.
\end{matrix}
$$
We will adjust our choices a little bit as follows.
Choose lifts $\tilde y_i \in M$ of $y_i \in I_\varphi$
and $\tilde s_i \in M$ of $s_i \in I_\psi$. It may not be the case
that $\mathfrak q \tilde y_1 \subset \langle x_1, \ldots, x_b\rangle$
and it may not be the case that
$\mathfrak q \tilde s_1 \subset \langle t_1, \ldots, t_a\rangle$.
However, using that $\mathfrak q$ is finitely generated (as in the proof
of Lemma \ref{lemma-good-sequence-exists}) we can find a
$d \in R$, $d \not \in \mathfrak q$ such that
$\mathfrak q d\tilde y_1 \subset \langle x_1, \ldots, x_b\rangle$
and
$\mathfrak q d\tilde s_1 \subset \langle t_1, \ldots, t_a\rangle$.
Thus after replacing $y_i$ by $dy_i$,
$\tilde y_i$ by $d\tilde y_i$, $s_i$ by $ds_i$ and $\tilde s_i$
by $d\tilde s_i$ we see that we may assume also that
$x_1, \ldots, x_b, \tilde y_1, \ldots, \tilde y_b$
and $t_1, \ldots, t_a, \tilde s_1, \ldots, \tilde s_b$
are good sequences in $M$.

\medskip\noindent
Finally, we choose a good sequence
$z_1, \ldots, z_l$ in the finite $R$-module
$$
\langle
x_1, \ldots, x_b, \tilde y_1, \ldots, \tilde y_a
\rangle
\cap
\langle
t_1, \ldots, t_a, \tilde s_1, \ldots, \tilde s_b
\rangle.
$$
Note that this is also a good sequence in $M$.

\medskip\noindent
Since $I_{\varphi, \mathfrak q} = K_{\psi, \mathfrak q}$
there is a unique element $h \in \kappa(\mathfrak q)$ such that
$[y_1, \ldots, y_a] = h [t_1, \ldots, t_a]$
inside $\det_{\kappa(\mathfrak q)}(K_{\psi, \mathfrak q})$.
Similarly, as $I_{\psi, \mathfrak q} = K_{\varphi, \mathfrak q}$
there is a unique element $h \in \kappa(\mathfrak q)$ such that
$[s_1, \ldots, s_b] = g [x_1, \ldots, x_b]$
inside $\det_{\kappa(\mathfrak q)}(K_{\varphi, \mathfrak q})$.
We can also do this with the three good sequences we have
in $M$. All in all we get the following identities
\begin{align*}
[y_1, \ldots, y_a]
& =
h [t_1, \ldots, t_a] \\
[s_1, \ldots, s_b]
& =
g [x_1, \ldots, x_b] \\
[z_1, \ldots, z_l]
& =
f_\varphi [x_1, \ldots, x_b, \tilde y_1, \ldots, \tilde y_a] \\
[z_1, \ldots, z_l]
& =
f_\psi [t_1, \ldots, t_a, \tilde s_1, \ldots, \tilde s_b]
\end{align*}
for some $g, h, f_\varphi, f_\psi \in \kappa(\mathfrak q)$.

\medskip\noindent
Having set up all this
notation let us compute $\det_{\kappa(\mathfrak q)}(M, \varphi, \psi)$.
Namely, consider the element $[z_1, \ldots, z_l]$.
Under the map $\gamma_\psi \circ \sigma \circ \gamma_\varphi^{-1}$
of Definition \ref{definition-periodic-determinant} we have
\begin{eqnarray*}
[z_1, \ldots, z_l] & = &
f_\varphi [x_1, \ldots, x_b, \tilde y_1, \ldots, \tilde y_a] \\
& \mapsto & f_\varphi [x_1, \ldots, x_b] \otimes [y_1, \ldots, y_a] \\
& \mapsto &
f_\varphi h/g [t_1, \ldots, t_a] \otimes [s_1, \ldots, s_b] \\
& \mapsto &
f_\varphi h/g [t_1, \ldots, t_a, \tilde s_1, \ldots, \tilde s_b] \\
& = &
f_\varphi h/f_\psi g [z_1, \ldots, z_l]
\end{eqnarray*}
This means that
$\det_{\kappa(\mathfrak q)}
(M_{\mathfrak q}, \varphi_{\mathfrak q}, \psi_{\mathfrak q})$
is equal to $f_\varphi h/f_\psi g$ up to a sign.

\medskip\noindent
We abbreviate the following quantities
\begin{eqnarray*}
k_\varphi & = & \text{length}_R(K_\varphi/\langle x_1, \ldots, x_b\rangle) \\
k_\psi    & = & \text{length}_R(K_\psi/\langle t_1, \ldots, t_a\rangle) \\
i_\varphi & = & \text{length}_R(I_\varphi/\langle y_1, \ldots, y_a\rangle) \\
i_\psi    & = & \text{length}_R(I_\psi/\langle s_1, \ldots, s_a\rangle) \\
m_\varphi & = & \text{length}_R(M/
\langle x_1, \ldots, x_b, \tilde y_1, \ldots, \tilde y_a\rangle) \\
m_\psi    & = & \text{length}_R(M/
\langle t_1, \ldots, t_a, \tilde s_1, \ldots, \tilde s_b\rangle) \\
\delta_\varphi & = & \text{length}_R(
\langle x_1, \ldots, x_b, \tilde y_1, \ldots, \tilde y_a\rangle
\langle z_1, \ldots, z_l\rangle) \\
\delta_\psi & = & \text{length}_R(
\langle t_1, \ldots, t_a, \tilde s_1, \ldots, \tilde s_b\rangle
\langle z_1, \ldots, z_l\rangle)
\end{eqnarray*}
Using the exact sequences $0 \to K_\varphi \to M \to I_\varphi \to 0$
we get $m_\varphi = k_\varphi + i_\varphi$. Similarly we have
$m_\psi = k_\psi + i_\psi$. We have
$\delta_\varphi + m_\varphi = \delta_\psi + m_\psi$ since this
is equal to the colength of $\langle z_1, \ldots, z_l \rangle$
in $M$. Finally, we have
$$
\delta_\varphi = \text{ord}_{R/\mathfrak q}(f_\varphi),
\quad
\delta_\psi = \text{ord}_{R/\mathfrak q}(f_\psi)
$$
by our first application of the key Lemma \ref{lemma-key-lemma}.

\medskip\noindent
Next, let us compute the multiplicity of the periodic complex
\begin{eqnarray*}
e_R(M, \varphi, \psi) & = &
\text{length}_R(K_\varphi/I_\psi) - \text{length}_R(K_\psi/I_\varphi) \\
& = &
\text{length}_R(
\langle x_1, \ldots, x_b\rangle/
\langle s_1, \ldots, s_b\rangle)
+ k_\varphi - i_\psi \\
& & -
\text{length}_R(
\langle t_1, \ldots, t_a\rangle/
\langle y_1, \ldots, y_a\rangle)
- k_\psi + i_\varphi \\
& = &
\text{ord}_{R/\mathfrak q}(g/h) + k_\varphi - i_\psi - k_\psi + i_\varphi \\
& = &
\text{ord}_{R/\mathfrak q}(g/h) + m_\varphi - m_\psi \\
& = &
\text{ord}_{R/\mathfrak q}(g/h) + \delta_\psi - \delta_\varphi \\
& = &
\text{ord}_{R/\mathfrak q}(f_\psi g/f_\varphi h)
\end{eqnarray*}
where we used the key Lemma \ref{lemma-key-lemma} twice in the third equality.
By our computation of $\det_{\kappa(\mathfrak q)}
(M_{\mathfrak q}, \varphi_{\mathfrak q}, \psi_{\mathfrak q})$
this proves the proposition.
\end{proof}
